\section{Predictions and a research agenda}\label{sec:agenda}
The argument yields predictions that could fail.
\begin{enumerate}
\item \textbf{P1 (skill-level moderation).} Baseline executive function predicts acquisition of restructuring skill more than acquisition of exposure-based skills, and predicts overall outcome only weakly.
\item \textbf{P2 (stress-dependence).} The gap between calm and in-the-moment reappraisal is larger in low-executive-function samples and at higher arousal than the moderate stressors used so far.
\item \textbf{P3 (component).} In adults with ADHD and anxiety, an exposure-and-behaviour protocol is not inferior to the same protocol plus restructuring.
\item \textbf{P4 (ADHD).} ADHD status predicts less benefit from the cognitive module specifically.
\end{enumerate}
\textbf{What would refute the argument:} an adequately powered study finding no moderation of skill acquisition by executive function; a component trial showing that restructuring adds substantial benefit over exposure; or meta-regression showing executive-function-heavy samples respond as well. The cheapest next steps are (i) an individual-participant or meta-regression analysis of existing trials (see \texttt{experiments/}), (ii) a PRISMA-style umbrella search to lift this review's method to the target journal's standard, and (iii) a power calculation for a moderator trial.

\section{Conclusion}
Restructuring-heavy CBT is effective on average, though less so against active controls and after bias adjustment, and its effects have not grown over three decades. What is not established is that the cognitive component is the active ingredient, or that it fails for people with low executive function. The best-supported claims are narrow: reappraisal draws on working memory; component studies, though underpowered, repeatedly fail to find an added benefit of restructuring over behavioural elements; executive capacity may limit the uptake of the cognitive skill but not overall response; and in adult ADHD the effective ingredients look behavioural. The strong, popular form of the argument, that anxiety takes the prefrontal cortex ``off-line,'' is not supported by the primary literature. Direct tests in adults with ADHD and anxiety are the missing evidence.
